\chapter{The Geometry That Moves Without Forgetting}

A quantum state is usually written with complex numbers. Split those numbers into real and imaginary parts, however, and the same evolution becomes a trajectory through a real phase space. The trajectory obeys Hamiltonian rules: it moves through paired coordinates, much like position and momentum.

That observation inspired the idea of a \emph{symplectic computer}. But a larger family of mathematical transformations does not automatically mean a more powerful computer. A machine must still tell us how states are prepared, how answers are measured, how much precision is required, and how much energy amplification costs.

For MMALS, the most useful lesson is different. Symplectic geometry can describe \textbf{how information and resources move} through an adaptive system.

\section*{Hosts, network, and exchange}
In MMALS, the hosts are the competency-bearing organisms. The mycelial network is the medium that transports signals and resources. The mycorrhizal interfaces are the exchange relationships between each host and that medium. Keeping these roles separate matters: the network does not magically own every competence, and a host does not coordinate the entire system alone.

Each host receives a paired state
\[
z_h=(q_h,p_h).
\]
One part can represent its current functional state; the other can represent momentum, pressure to adapt, or resource flow.

\section*{Movement is not learning}
Pure symplectic motion is reversible. It can circulate, rotate, and redistribute, but it cannot easily settle into a stable memory. MMALS therefore adds dissipation:
\[
\dot Z=[J-R]\nabla\mathcal H+Bu.
\]
Here $J$ transports and explores, $R$ settles and consolidates, and $Bu$ steers the system toward the current goal. A second geometry, Fisher--Rao geometry, updates beliefs about context, routes, competence, and uncertainty.

\section*{The honest claim}
This chapter does not claim that MMALS has become a quantum computer, or that symplectic computers have been proved superior to quantum machines. The claim is more practical and more demanding:

\begin{quote}
MMALS now has a candidate mathematical basement in which transport, learning, memory, specialization, resource allocation, and control can be expressed within one auditable dynamical system.
\end{quote}

The next step is experimental. Geometry is valuable only when it buys measurable capability under equal precision, energy, and memory budgets.
